\documentclass[../../../include/open-logic-section]{subfiles}

\begin{document}

\olfileid{sth}{ordinals}{vn}
\olsection[بناء فون نويمان]{بناء فون نويمان للأعداد الترتيبية}

تفتح~\olref[sth][ordinals][iso]{thm:woalwayscomparable} بابًا آخر:
لِمَ لا نجعل للترتيبات الجيدة أشياء تقوم مقامها، نسميها \emph{أنماط
الترتيب}؟ فإذا رمزنا بـ$\ordtype{A, <}$ إلى نمط ترتيب
$\tuple{A, <}$، كان المطلوب أن يتحقق المبدآن الآتيان:
\begin{align*}
\ordtype{A, <} = \ordtype{B, \lessdot} &
\text{ إذا وفقط إذا كان } \ordeq{\tuple{A, <}}{\tuple{B, \lessdot}}\\
\ordtype{A, <} < \ordtype{B, \lessdot}&
\text{ إذا وفقط إذا كان }\ordeq{\tuple{A, <}}{\tuple{B_b, \lessdot_b}}
\text{ لبعض }b \in B
\end{align*}
ونروم فوق ذلك أن تكون أنماط الترتيب \emph{مجموعات معينة}، على نحو
ما بنينا الأعداد الطبيعية من مجموعات معينة.

وأشهر السبل إلى هذا، وهو الذي نسلكه، تعريف الأنماط بمجموعات
\emph{قياسية} مرتبة ترتيبًا جيدًا. وأول من أدخل هذه المجموعات
فون نويمان.

\begin{defn}
المجموعة~$A$ \emph{متعدية} إذا وفقط إذا كان
$(\forall x \in A)x \subseteq A$. وتسمى~$A$ \emph{عددًا ترتيبيًا}
إذا وفقط إذا كانت~$A$ متعدية وكانت~$\in$ ترتبها ترتيبًا جيدًا.
\end{defn}
\noindent
ونرمز إلى الأعداد الترتيبية فيما يلي بحروف يونانية. فإذا كان
$\alpha$ عددًا ترتيبيًا، لزم من التعريف أن~$\tuple{\alpha, \in_\alpha}$
ترتيب جيد، حيث~$\in_\alpha = \Setabs{\tuple{x, y} \in \alpha^2}{x \in y}$.
ومن ثم نتسامح قليلًا في الترميز، فنقول إن~$\alpha$ \emph{نفسه}
ترتيب جيد.

وأوائل الأعداد الترتيبية هي:
\[
\emptyset, \{\emptyset\},
\{\emptyset, \{\emptyset\}\},
\{\emptyset, \{\emptyset\}, \{\emptyset, \{\emptyset\}\}\}, \ldots
\]
وهذه بعينها الأوائل التي مرت في مسلّمة اللانهاية، أعني في تعريف
فون نويمان لـ$\omega$؛ انظر~\olref[sth][z][infinity-again]{sec}.
ولا اتفاق عارض هنا: فتعريفه يجعل الأعداد الطبيعية أعدادًا ترتيبية،
ولا يقصر الأعداد الترتيبية عليها، بل يجيز ما تجاوز المتناهي أيضًا.

ويبقى السؤال المعهود: أهذه \emph{هي} الأعداد الترتيبية، أم إن فون
نويمان لم يعطنا إلا مجموعات يصح أن \emph{نعاملها} معاملتها؟ القول
في هذا نظير ما تقدم في~\olref[sfr][rel][ref]{sec}،
و\olref[sfr][arith][ref]{sec}، و\olref[sfr][infinite][dedekindsproof]{sec}،
و\olref[sth][z][nat]{sec}؛ فلا نعيد تفصيله. وإنما نعني من الآن
بـ«الأعداد الترتيبية» أعداد فون نويمان الترتيبية.

\end{document}
